\section*{Hoofdstuk \ref{ch:b3:environmental-toxicology} --- Milieuchemie en toxicologie}

\begin{solution}{exo:b3:environmental-toxicology:1}
Een homonucleair diatomisch molecuul heeft één vibratie, die het
dipoolmoment nul houdt, en argon heeft geen vibratie: geen van beide kan
infrarode straling absorberen. $\ce{CO2}$ absorbeert via
zijn buigvibratie (en zijn antisymmetrische strekvibratie), die een
oscillerende dipool scheppen; de buigvibratie ligt midden in de emissie
van de aarde.
\end{solution}

\begin{solution}{exo:b3:environmental-toxicology:2}
$10^{0.1} = 1.26$: $[\ce{H+}]$ stijgt met 26\,\%.
\end{solution}

\begin{solution}{exo:b3:environmental-toxicology:3}
$(2.0 + 0.50)\ \unit{mmol/L} \times \qty{100.1}{mg/mmol} = \qty{250}{mg/L}$ $\ce{CaCO3}$.
\end{solution}

\begin{solution}{exo:b3:environmental-toxicology:4}
N: $0.60/14.0 = \qty{0.043}{mmol/L}$; P: $0.010/31.0 = \qty{3.2e-4}{mmol/L}$; N\,:\,P $= 133$, ver boven 16: de fosfor raakt
eerst op en beperkt de groei.
\end{solution}

\begin{solution}{exo:b3:environmental-toxicology:5}
$L_0 = \mathrm{BZV_5}/(1 - \eu^{-0.23 \times 5}) = 170/0.683 = \qty{250}{mg/L}$.
\end{solution}

\begin{solution}{exo:b3:environmental-toxicology:6}
$\qty{4.40e-3}{L} \times \qty{0.0200}{mol/L} = \qty{8.80e-5}{mol}$ in \qty{0.1000}{L}: \omterm{def:b3:environmental-toxicology:alkalinity}{alkaliniteit} $\qty{8.80e-4}{mol/L}$; als $\ce{HCO3-}$ (\qty{61.0}{g/mol})
\qty{53.7}{mg/L}.
\end{solution}

\begin{solution}{exo:b3:environmental-toxicology:7}
$K_d = 0.020 \times 300 = \qty{6.0}{L/kg}$. Gesorbeerd/totaal $= K_dm_s/(K_dm_s + V_w) = 9.0/(9.0 + 0.30) = 0.97$: 97\,\% gesorbeerd.
\end{solution}

\begin{solution}{exo:b3:environmental-toxicology:8}
Benzeen: $\log\mathrm{BCF} = 0.85 \times 2.13 - 0.70 = 1.11$, BCF ongeveer 13. PCB: $0.85 \times 6.67 - 0.70 = 4.97$, BCF ongeveer
\num{9e4}: het concentreert zich in vissen zo'n zevenduizend keer sterker
dan benzeen; bij zo'n
$K_{ow}$ zit de regressie dicht bij de grens van haar geldigheid.
\end{solution}

\begin{solution}{exo:b3:environmental-toxicology:9}
$\qty{192}{mg/kg} \times \qty{70}{kg} = \qty{13}{g}$, ongeveer 150 koppen. Soorten verschillen in opname en
metabolisme,
en een $\mathrm{LD_{50}}$ is geen veilige drempel; schadelijke effecten beginnen ver daaronder.
\end{solution}

\begin{solution}{exo:b3:environmental-toxicology:10}
PNEC $= \qty{0.50}{mg/L}/1000 = \qty{0.50}{\micro g/L}$; $\mathrm{RQ} = 0.20/0.50 = 0.40$, onder 1: geen zorg bij deze blootstelling.
\end{solution}

\begin{solution}{exo:b3:environmental-toxicology:11}
Plankton/water: $\qty{0.015}{mg/kg}/\qty{2.0e-6}{mg/L} = \qty{7.5e3}{L/kg}$ (bioconcentratie). Vis/plankton: $0.30/0.015 = 20$;
vogel/vis: $6.0/0.30 = 20$ (\omterm{def:b3:environmental-toxicology:bio}{biomagnificatie} bij elke stap); vogel/water: $\num{3e6}$.
\end{solution}

\begin{solution}{exo:b3:environmental-toxicology:12}
$\qty{1.0e9}{g}/\qty{64.1}{g/mol} = \qty{1.56e7}{mol}$ $\ce{SO2}$, wat evenveel mol $\ce{H2SO4}$ geeft: $\qty{1.53e9}{g}$, ongeveer \qty{1500}{t}. Het
zet $\qty{3.12e7}{mol}$ $\ce{H+}$ vrij; bij pH 4.0 is $[\ce{H+}] = \qty{1.0e-4}{mol/L}$: $\qty{3.1e11}{L}$ regen.
\end{solution}

\begin{solution}{pb:b3:environmental-toxicology:1}
\textbf{1.} Het is tienduizend keer beter oplosbaar in octanol dan in
water: hydrofoob, het zal op organische stof sorberen en zich in vet
ophopen.

\textbf{2.} $K_{oc} = 0.41 \times 10^4 = \qty{4100}{L/kg}$; $K_{sw} = f_{oc}K_{oc} = 0.040 \times 4100 = \qty{164}{L/kg}$.

\textbf{3.} $K_{fw} = 0.048 \times 10^4 = \qty{480}{L/kg}$.

\textbf{4.} Nee: in evenwicht bevat de lucht $10^{-4}$ van de waterconcentratie.

\textbf{5.} Het neutrale, hydrofobe molecuul lost op in de organische
stof, zoals in octanol; minerale oppervlakken zijn polair en met water
bedekt.

\textbf{6.} Het sediment, waarvan de capaciteit, $K_{sw}m_s$, het grootst is.

\textbf{7.} $V_w = \qty{1.0e9}{L}$; $K_{aw}V_a = \qty{1.0e8}{L}$; $K_{sw}m_s = \qty{1.64e10}{L}$; $K_{fw}m_f = \qty{4.8e6}{L}$.

\textbf{8.} $C_w = \qty{1.0e8}{mg}/\qty{1.75e10}{L} = \qty{5.7e-3}{mg/L}$, \qty{5.7}{\micro g/L}.

\textbf{9.} Water \qty{5.7}{kg}, lucht \qty{0.57}{kg}, sediment \qty{94}{kg}, vis \qty{0.027}{kg}.

\textbf{10.} Lucht $\qty{5.7e-7}{mg/L}$; sediment \qty{0.94}{mg/kg}; vis \qty{2.7}{mg/kg}.

\textbf{11.} Evenwicht tussen alle compartimenten tegelijk, zonder afbraak,
zonder in- of uitstroom, en met goed gemengde compartimenten.

\textbf{12.} Met afbraak en stromen in stationaire toestand (niveau II)
wordt de inventaris bepaald door de snelheden van aanvoer en verlies; met
overdrachtssnelheden tussen de compartimenten (niveau III) staan de
concentraties niet meer in evenwichtsverhoudingen, en wordt het
compartiment dat de aanvoer ontvangt verrijkt.

\textbf{13.} $k = \ln 2/\qty{60}{d} = \qty{0.0116}{d^{-1}}$.

\textbf{14.} $\eu^{-0.0116 \times 365} = 0.0145$: 1.5\,\% blijft over.

\textbf{15.} $\qty{5.7}{\micro g/L} \times 0.0145 = \qty{0.083}{\micro g/L}$.

\textbf{16.} De afbraak is in koud, zuurstofloos sediment vaak trager, en
het gesorbeerde pesticide komt traag weer vrij in het water, zodat het
aanwezig blijft nadat het water al gezuiverd is.

\textbf{17.} Bindingen die hydrolyse, oxidatie en microbiële aanval
weerstaan (aromatische C--Cl bijvoorbeeld), een lage oplosbaarheid in
water, en sorptie die het tegen afbrekende organismen beschermt.

\textbf{18.} PNEC $= \qty{50}{\micro g/L}/10 = \qty{5.0}{\micro g/L}$.

\textbf{19.} $\mathrm{RQ} = 5.7/5.0 = 1.1$.

\textbf{20.} Net boven 1: er wordt een risico aangegeven, en de
beoordeling moet verfijnd worden (betere toxiciteitsgegevens, gemeten
concentraties) of de blootstelling verminderd.

\textbf{21.} \qty{2.7}{mg/kg}, bijna drie keer de drempel van \qty{1.0}{mg/kg}: de vis mag niet gegeten worden
tot de concentratie gedaald is.

\textbf{22.} Concentraties in water, sediment en vis over de tijd en op
meerdere punten, om het model te toetsen en de daling te volgen.

\textbf{23.} Minder pesticide toepassen, toepassing weg van regen en van
de oever, bufferstroken die afspoeling tegenhouden, of een minder
persistent alternatief.

\textbf{24.} Ze deelt het gemeten niveau zonder effect om de kloof te
overbruggen tussen enkele laboratoriumsoorten en een heel ecosysteem;
een factor 10 in plaats van 1000 verandert de PNEC, en het oordeel,
honderdvoudig.

\textbf{25.} Het \omterm{def:b3:environmental-toxicology:ecotoxicology}{risicoquotiënt} PEC/PNEC is ongeveer 1.1, net boven 1.
\end{solution}
